\chapter{Objetivos}

\section{Objetivo general}

Dise\~nar y desarrollar un prototipo funcional de software que analice resonancias magn\'eticas lumbares sagitales para segmentar estructuras anat\'omicas relevantes, obtener un conjunto acotado de mediciones cuantitativas y generar una salida visual y estructurada editable que asista al profesional sin reemplazar su criterio cl\'inico.

\section{Objetivos espec\'ificos}

Para alcanzar el objetivo general, se definen los siguientes objetivos espec\'ificos:

\begin{enumerate}
    \item Relevar antecedentes acad\'emicos, datasets p\'ublicos, herramientas abiertas y soluciones comerciales relacionadas con el an\'alisis asistido de RM lumbar.
    \item Identificar la brecha funcional que justifica un prototipo acad\'emico orientado a segmentaci\'on, mediciones, visualizaci\'on y salida editable.
    \item Definir el alcance del MVP, incluyendo estructuras anat\'omicas, funcionalidades, validaci\'on prevista y exclusiones cl\'inicas, t\'ecnicas y regulatorias.
    \item Analizar necesidades de potenciales usuarios mediante actividades de user research vinculadas con la revisi\'on de RM lumbar.
    \item Dise\~nar e implementar un pipeline de carga, normalizaci\'on, preprocesamiento, segmentaci\'on, postprocesamiento y visualizaci\'on.
    \item Adaptar o implementar un modelo de segmentaci\'on para v\'ertebras, discos intervertebrales y canal espinal en RM lumbar sagital.
    \item Desarrollar mediciones geom\'etricas simples derivadas de las m\'ascaras generadas.
    \item Dise\~nar una interfaz y una salida estructurada editable que permitan revisar resultados sin constituir un informe cl\'inico autom\'atico.
    \item Evaluar t\'ecnicamente el prototipo mediante m\'etricas de segmentaci\'on y comparaci\'on contra anotaciones de referencia.
    \item Complementar la validaci\'on t\'ecnica con revisi\'on cualitativa profesional y documentar decisiones, limitaciones, resultados y trabajos futuros.
\end{enumerate}
